% !TEX root = ../thesis.tex
% ==============================================================================
%  ACKNOWLEDGMENTS PAGE (frontmatter/acknowledgments.tex)
% ==============================================================================
%
%  This file contains the acknowledgments section (صفحه سپاس و قدردانی).
%
% ------------------------------------------------------------------------------
%  HOW TO USE & CUSTOMIZE:
% ------------------------------------------------------------------------------
%  - Wrap your text inside:
%    \begin{thesisacknowledgments} ... \end{thesisacknowledgments}.
%  - The environment automatically sets up the centered header and spacing.
%  - To adjust the vertical spacing between the title and text, customize
%    \setlength{\thesisackbeforespace}{...} in `styles/custom.tex`.
%  - To toggle this page in the thesis:
%    Toggle `% !TEX root = ../thesis.tex
% ==============================================================================
%  ACKNOWLEDGMENTS PAGE (frontmatter/acknowledgments.tex)
% ==============================================================================
%
%  This file contains the acknowledgments section (صفحه سپاس و قدردانی).
%
% ------------------------------------------------------------------------------
%  HOW TO USE & CUSTOMIZE:
% ------------------------------------------------------------------------------
%  - Wrap your text inside:
%    \begin{thesisacknowledgments} ... \end{thesisacknowledgments}.
%  - The environment automatically sets up the centered header and spacing.
%  - To adjust the vertical spacing between the title and text, customize
%    \setlength{\thesisackbeforespace}{...} in `styles/custom.tex`.
%  - To toggle this page in the thesis:
%    Toggle `% !TEX root = ../thesis.tex
% ==============================================================================
%  ACKNOWLEDGMENTS PAGE (frontmatter/acknowledgments.tex)
% ==============================================================================
%
%  This file contains the acknowledgments section (صفحه سپاس و قدردانی).
%
% ------------------------------------------------------------------------------
%  HOW TO USE & CUSTOMIZE:
% ------------------------------------------------------------------------------
%  - Wrap your text inside:
%    \begin{thesisacknowledgments} ... \end{thesisacknowledgments}.
%  - The environment automatically sets up the centered header and spacing.
%  - To adjust the vertical spacing between the title and text, customize
%    \setlength{\thesisackbeforespace}{...} in `styles/custom.tex`.
%  - To toggle this page in the thesis:
%    Toggle `% !TEX root = ../thesis.tex
% ==============================================================================
%  ACKNOWLEDGMENTS PAGE (frontmatter/acknowledgments.tex)
% ==============================================================================
%
%  This file contains the acknowledgments section (صفحه سپاس و قدردانی).
%
% ------------------------------------------------------------------------------
%  HOW TO USE & CUSTOMIZE:
% ------------------------------------------------------------------------------
%  - Wrap your text inside:
%    \begin{thesisacknowledgments} ... \end{thesisacknowledgments}.
%  - The environment automatically sets up the centered header and spacing.
%  - To adjust the vertical spacing between the title and text, customize
%    \setlength{\thesisackbeforespace}{...} in `styles/custom.tex`.
%  - To toggle this page in the thesis:
%    Toggle `\include{frontmatter/acknowledgments}` in `thesis.tex`.
% ==============================================================================

\begin{thesisacknowledgments}
این قالب حروف‌چینی بر پایه‌ی قالب اولیه‌ی پایان‌نامه دانشگاه صنعتی شریف بنا شده است که توسط جناب آقای دکتر حمید ضرابی‌زاده طراحی گردید و با مشارکت آقایان احسان امام‌جمعه‌زاده و امید غیبی تکمیل شد. این اثر نقشی بنیادین در استانداردسازی حروف‌چینی پایان‌نامه‌های دانشگاهی با استفاده از زی‌پرشین داشته است.

نسخه‌ی بازطراحی‌شده‌ی حاضر، توسط \textbf{مریم رضائی} با هدف مدرن‌سازی زیرساخت فنی، ارتقای سهولت کاربری، و بهبود طراحی ظاهری پایان‌نامه بازنویسی شده است. محورهای اصلی این بازنگری شامل تفکیک ساختاری هسته‌ی حروف‌چینی از لایه‌ی تنظیمات کاربر، طراحی سامانه‌ی مرتب‌سازی خودکار واژه‌نامه‌های دوزبانه بر اساس ترتیب الفبای فارسی، اتوماتیک کردن جزئیات فرم‌های اداری تصویب‌نامه و اظهارنامه به همراه اسکریپت خروجی‌گیر مستقل، ارتقای سبک مراجع‌نگاری به بسته‌ی طراحی شده‌ی \code{unsrtabbrv-modern.bst} با پیوندهای برخط و سایر استانداردسازی‌ها، ادغام خانواده قلم ادوبی کَسِلون برای حروف‌چینی لاتین، و رفع کامل هشدارهای کامپایلر می‌باشد.
\end{thesisacknowledgments}
` in `thesis.tex`.
% ==============================================================================

\begin{thesisacknowledgments}
این قالب حروف‌چینی بر پایه‌ی قالب اولیه‌ی پایان‌نامه دانشگاه صنعتی شریف بنا شده است که توسط جناب آقای دکتر حمید ضرابی‌زاده طراحی گردید و با مشارکت آقایان احسان امام‌جمعه‌زاده و امید غیبی تکمیل شد. این اثر نقشی بنیادین در استانداردسازی حروف‌چینی پایان‌نامه‌های دانشگاهی با استفاده از زی‌پرشین داشته است.

نسخه‌ی بازطراحی‌شده‌ی حاضر، توسط \textbf{مریم رضائی} با هدف مدرن‌سازی زیرساخت فنی، ارتقای سهولت کاربری، و بهبود طراحی ظاهری پایان‌نامه بازنویسی شده است. محورهای اصلی این بازنگری شامل تفکیک ساختاری هسته‌ی حروف‌چینی از لایه‌ی تنظیمات کاربر، طراحی سامانه‌ی مرتب‌سازی خودکار واژه‌نامه‌های دوزبانه بر اساس ترتیب الفبای فارسی، اتوماتیک کردن جزئیات فرم‌های اداری تصویب‌نامه و اظهارنامه به همراه اسکریپت خروجی‌گیر مستقل، ارتقای سبک مراجع‌نگاری به بسته‌ی طراحی شده‌ی \code{unsrtabbrv-modern.bst} با پیوندهای برخط و سایر استانداردسازی‌ها، ادغام خانواده قلم ادوبی کَسِلون برای حروف‌چینی لاتین، و رفع کامل هشدارهای کامپایلر می‌باشد.
\end{thesisacknowledgments}
` in `thesis.tex`.
% ==============================================================================

\begin{thesisacknowledgments}
این قالب حروف‌چینی بر پایه‌ی قالب اولیه‌ی پایان‌نامه دانشگاه صنعتی شریف بنا شده است که توسط جناب آقای دکتر حمید ضرابی‌زاده طراحی گردید و با مشارکت آقایان احسان امام‌جمعه‌زاده و امید غیبی تکمیل شد. این اثر نقشی بنیادین در استانداردسازی حروف‌چینی پایان‌نامه‌های دانشگاهی با استفاده از زی‌پرشین داشته است.

نسخه‌ی بازطراحی‌شده‌ی حاضر، توسط \textbf{مریم رضائی} با هدف مدرن‌سازی زیرساخت فنی، ارتقای سهولت کاربری، و بهبود طراحی ظاهری پایان‌نامه بازنویسی شده است. محورهای اصلی این بازنگری شامل تفکیک ساختاری هسته‌ی حروف‌چینی از لایه‌ی تنظیمات کاربر، طراحی سامانه‌ی مرتب‌سازی خودکار واژه‌نامه‌های دوزبانه بر اساس ترتیب الفبای فارسی، اتوماتیک کردن جزئیات فرم‌های اداری تصویب‌نامه و اظهارنامه به همراه اسکریپت خروجی‌گیر مستقل، ارتقای سبک مراجع‌نگاری به بسته‌ی طراحی شده‌ی \code{unsrtabbrv-modern.bst} با پیوندهای برخط و سایر استانداردسازی‌ها، ادغام خانواده قلم ادوبی کَسِلون برای حروف‌چینی لاتین، و رفع کامل هشدارهای کامپایلر می‌باشد.
\end{thesisacknowledgments}
` in `thesis.tex`.
% ==============================================================================

\begin{thesisacknowledgments}
این قالب حروف‌چینی بر پایه‌ی قالب اولیه‌ی پایان‌نامه دانشگاه صنعتی شریف بنا شده است که توسط جناب آقای دکتر حمید ضرابی‌زاده طراحی گردید و با مشارکت آقایان احسان امام‌جمعه‌زاده و امید غیبی تکمیل شد. این اثر نقشی بنیادین در استانداردسازی حروف‌چینی پایان‌نامه‌های دانشگاهی با استفاده از زی‌پرشین داشته است.

نسخه‌ی بازطراحی‌شده‌ی حاضر، توسط \textbf{مریم رضائی} با هدف مدرن‌سازی زیرساخت فنی، ارتقای سهولت کاربری، و بهبود طراحی ظاهری پایان‌نامه بازنویسی شده است. محورهای اصلی این بازنگری شامل تفکیک ساختاری هسته‌ی حروف‌چینی از لایه‌ی تنظیمات کاربر، طراحی سامانه‌ی مرتب‌سازی خودکار واژه‌نامه‌های دوزبانه بر اساس ترتیب الفبای فارسی، اتوماتیک کردن جزئیات فرم‌های اداری تصویب‌نامه و اظهارنامه به همراه اسکریپت خروجی‌گیر مستقل، ارتقای سبک مراجع‌نگاری به بسته‌ی طراحی شده‌ی \code{unsrtabbrv-modern.bst} با پیوندهای برخط و سایر استانداردسازی‌ها، ادغام خانواده قلم ادوبی کَسِلون برای حروف‌چینی لاتین، و رفع کامل هشدارهای کامپایلر می‌باشد.
\end{thesisacknowledgments}
